\documentclass[11pt,a4paper]{ctexart}

\usepackage[a4paper,left=3cm,right=3cm,top=2.8cm,bottom=2.8cm]{geometry}
\usepackage{fontspec}
\usepackage{xeCJK}
\usepackage{setspace}
\usepackage{amsmath,amssymb}
\usepackage{booktabs}
\usepackage{array}
\usepackage{tabularx}
\usepackage{hyperref}

\setmainfont{DejaVu Serif}
\setsansfont{DejaVu Sans}
\setmonofont{DejaVu Sans Mono}[Scale=0.88]
\setCJKmainfont{Noto Serif SC}[AutoFakeBold=2.5,AutoFakeSlant=0.15]
\setCJKsansfont{Noto Sans SC}[AutoFakeBold=2.5]
\setCJKmonofont{Noto Sans SC}

\onehalfspacing
\setlength{\emergencystretch}{1em}

\ctexset{
  section = {
    format = \Large\sffamily\bfseries\raggedright,
    aftername = \hspace{0.5em},
  },
  subsection = {
    format = \large\sffamily\bfseries\raggedright,
    aftername = \hspace{0.5em},
  }
}

\hypersetup{
  colorlinks=true,
  linkcolor=black,
  citecolor=black,
  urlcolor=black,
  pdftitle={论主等于六：基于阿奎那五路证明与 Lean 4 形式验证的第一算术本原考辨},
  pdfauthor={Chan Isidore}
}

\begin{document}
\pagestyle{plain}

\begin{center}
  {\LARGE\sffamily\bfseries 论主等于六：基于阿奎那五路证明与 Lean 4 形式验证的第一算术本原考辨}\\[1.2em]
  {\normalsize \textbf{作者}：Chan Isidore}\\[0.2em]
{\small 聖神修院（Holy Spirit Seminary），香港}\\[0.2em]
{\small \href{mailto:chengxi271@gmail.com}{chengxi271@gmail.com}}\\[0.4em]
  {\small \textbf{分类号}：MSC 03B35, 11A25, 00A30}
\end{center}

\vspace{0.8em}

\section*{摘要}

第一本原``主''在自然数秩序中居于何值，传统探讨多止步于定性思辨。本文依循阿奎那《神学大全》``五路证明''之体例，在交互式定理证明器 Lean 4 中将``主''设为待定自然数参数，考察其在变动止息、动力因构成、偶奇必然承载、完满等级居首与治理配额闭合五重规定下的算术归宿。形式验证表明：由真因子求和映射下的第一正不动点、真因子1与2与3对总体的穷尽构成、同时容纳2与3整除的最小正承载者、完全数序列的首要地位，以及二分与三分加一单位恰尽整体的配额秩序出发，五条路径各自推出主等于6，且与该定值互为充要条件。针对0、1、28与4等候选数的辩难进一步证实，唯有6统摄五路规定。

\noindent\textbf{关键词}：主等于六；五路证明；阿奎那；完全数；Lean 4；算术本原

\vspace{0.6em}
\hrule
\vspace{0.8em}

\section{引言}

探究万有秩序，首在追问第一本原（``主''）在自然数中是否有确定本位。旧说或谓其超逾定量，或轻推``一''``十''居尊；然若第一本原果真统摄生成、因果、必然、完满与治理，其在算术秩序中必落于唯一不亏不盈之定数。

阿奎那《神学大全》立``由果溯因''之五路以指认第一本原~\cite{aquinas}，惜止于自然哲学定性；奥古斯丁与欧几里得虽明``六''系首个完全数~\cite{augustine,euclid}，亦未以五路架构锁定其为``主''之唯一数值解，且常言推演易在零值或截断除法处混入伪解。

本文在 Lean 4 中将 $\text{主} : \mathbb{N}$ 设为待定参数，不预设常量亦不加自定义公理，依五路脉络考察其算术职分。形式验证表明，五路前提各自充要导出 $\text{主}=6$，并廓清 $0, 1, 4, 10, 12, 28$ 诸数之异议，遂定``主等于六''之论。

\section{相关工作}

\subsection{亚里士多德—阿奎那传统中的五路本原论}

以果溯因植根于亚里士多德对运动与原因链的考察：其《物理学》第八卷由受动链不可无限倒退确立第一不动者~\cite{physics}，《形而上学》第二卷阐明动力因与目的因同样拒绝无穷后退，且使他物成真者必自居最高完满~\cite{metaphysics}。阿奎那在《神学大全》第一部第二题中承此立``五路证明''，以``异议—答复''体式指认同一第一本原~\cite{aquinas}，成为本文的方法论轨范。

\subsection{完全数之教父神学诠释与古典算术构造}

关于完满性何以寄于``六''，奥古斯丁在《上帝之城》中指出，六日创世旨在以``六''表征工程之完满：它是首个由真部分（小于自身之正因子）精确合成的数，六分之一（$1$）、三分之一（$2$）与二分之一（$3$）之和恰为 $6$；反观十与十二之真部分和则为 $8$ 与 $16$~\cite{augustine}。

欧几里得《几何原本》给出偶完全数生成定理：自 $1$ 起按二倍比累取，若和为素数，则该素数乘末项即得完全数；最小序列 $1+2=3$ 乘 $2$ 得首数 $6$，次小由 $1+2+4=7$ 乘 $4$ 得 $28$~\cite{euclid}。故六兼具``真部分穷尽整体''与``完满居首''二性。

\subsection{依赖类型论与形式定理证明}

为免自然语言歧义，可将推演置于形式系统中检验。de Moura 等设计的 Lean 基于含归纳类型的构造演算并遵循小内核原则，一切底层证明项均须经内核类型检查~\cite{lean2015}；Lean 4 进而统一函数求值与形式演绎，支持对可计算列表函数直接施加量化证明~\cite{lean4}。依其定理证明规范，定义归约、可判定核验与线性算术皆立于标准公理之上~\cite{tpil}，使五路规约无需新增公理即可完成核验。

既有研究如 Benzmüller 等人在 Isabelle/HOL 高阶模态逻辑中检验 Gödel 论证并揭示特定版本之模态坍塌~\cite{modal}。本文不引入模态公理，径以自然数算术在 Lean 内核完成核验，零自定义公理。

\section{方法与形式规约}

本文在 Lean 4（版本 4.34.1，仅导入标准库 \texttt{Std}）中展开形式构造，源码见公开仓库 \texttt{lean/Main.lean}。对任意自然数 $n \in \mathbb{N}$，考察其小于自身的正整除因子。定义真部分列表 $P(n)$ 为区间 $[1, n-1]$ 中整除 $n$ 的升序自然数列表，真因子流溢映射 $s(n)$ 为 $n$ 的全体真部分之和：
\[
P(n) = [d \in \{1, \dots, n-1\} \mid n \bmod d = 0], \quad s(n) = \sum_{d \in P(n)} d
\]
当一自然数切实实存（$n > 0$），且在真因子流溢下既不减损（$s(n) < n$）亦不外溢（$s(n) > n$）时，称其为完全数，记作谓词 $K(n) \iff n > 0 \land s(n) = n$。

形式推演将第一本原声明为待定自然数参数 $\text{主} : \mathbb{N}$，不借常数赋值预先指定，而令候选数凭自身算术本性承受表~\ref{tab:five_ways} 所列五路显式规约之检验，全程不使用未证断言。

\section{结果}

\subsection{五路推演与算术本原之会归}

将待定参数 $\text{主}$ 置于上述五重规约之下，五条路径分别从变动、因果、必然、完满与治理出发，指向同一数值归宿（表~\ref{tab:five_ways}）。

\begin{table}[htbp]
  \centering
  \caption{五路证明之算术规定与共同归结}
  \label{tab:five_ways}
  \small
  \begin{tabularx}{\textwidth}{>{\raggedright\arraybackslash}p{3.4cm} >{\raggedright\arraybackslash}X c}
    \toprule
    路径 & 作用于待定参数 $\text{主}$ 的关键规定 & 归结 \\
    \midrule
    第一路：运动与不动者 & $\text{主}>0,\; s(\text{主})=\text{主},\; \forall m\in(0,\text{主}),\, s(m)\neq m$ & $\text{主} = 6$ \\
    第二路：动力因之总体 & $P(\text{主})=[1,2,3],\; s(\text{主})=\text{主}$ & $\text{主} = 6$ \\
    第三路：可能与必然性 & $\text{主}>0,\; 2\mid\text{主},\; 3\mid\text{主},\; \forall m>0,\, (2\mid m \land 3\mid m \to \text{主}\le m)$ & $\text{主} = 6$ \\
    第四路：完满性之等级 & $K(\text{主}),\; \forall m,\, K(m)\to \text{主}\le m$ & $\text{主} = 6$ \\
    第五路：秩序与目的治理 & $\text{主}>0,\; 2\mid\text{主},\; 3\mid\text{主},\; \text{主}/2+\text{主}/3+1=\text{主}$ & $\text{主} = 6$ \\
    \bottomrule
  \end{tabularx}
  \par\vspace{0.4em}
  {\footnotesize\raggedright 注：表中 $P(n)$ 为升序正真因子列表，$s(n)$ 为真因子和，$K(n)$ 表示完全数性质 $n>0 \land s(n)=n$。数据来源：本文定义与形式验证。\par}
\end{table}

\noindent\textbf{第一路：由真因子流溢之变动推证第一正不动者。}\\
凡自然数置于流溢映射 $n \mapsto s(n)$ 下，若其真部分之和异于自身（$s(n) \neq n$），则必处于由自身向他数迁流之变动中~\cite{aquinas,physics}。今逐一审视正整数之始：单子 $1$ 无真部分，一经流溢即坠入 $s(1)=0$；素数 $2, 3, 5$ 仅含真部分 $1$，流溢后皆退缩为 $1$；合数 $4$ 虽得真部分 $1$ 与 $2$，其和 $s(4)=3$ 仍不免于亏缺。可见自 $1$ 至 $5$ 诸数皆在流溢中失其本位，故任何满足 $s(\text{主})=\text{主}$ 的正不动者必满足 $\text{主} \ge 6$。反之，若 $\text{主} > 6$，则 $6$ 作为小于 $\text{主}$ 的正数，其真部分 $1, 2, 3$ 之和恰为 $s(6)=6$，早已在先实现了流溢之止息，这便与第一不动者``先于它的一切正数皆变动不居''的规定直接冲突。既不得小于 $6$，亦不得大于 $6$，故第一正不动者必为 $\text{主}=6$。

\medskip
\noindent\textbf{第二路：由原初动力因之总体推证果之完全生成。}\\
依第二路之理，任何复合事物皆不能充当自身的真部分（故 $n \notin P(n)$），其生成必仰赖先于自身的真因之合力~\cite{aquinas,metaphysics}。在正整数的生成序列中，最先显明自身特性的三项原初要素，乃是作为普遍起点的单子 $1$、开启偶分秩序的 $2$，以及开启奇素秩序的 $3$。当``主''以且仅以这三项原初要素为全部真部分（$P(\text{主})=[1, 2, 3]$），且三因之和不假外求、丝毫不差地穷尽果之总体（$s(\text{主})=\text{主}$）时，因果之链遂告闭合。将真部分列表直接代入流溢求和，立得 $\text{主} = 1 + 2 + 3 = 6$。

\medskip
\noindent\textbf{第三路：由偶奇分殊之可能性推证最小共同必然承载者。}\\
算术世界中，二元之偶性（以 $2$ 为本）与三元所表征之奇素本原（以首奇素数 $3$ 为本）本各自独立：仅具偶性者（如 $2, 4$）不能容纳三元之整除，仅具奇性者（如 $3$）亦不能容纳二元之平分。若无一共同的必然存在者同时整除承载二者（$2 \mid \text{主} \land 3 \mid \text{主}$），则偶与奇便终古隔绝而无从会通。在区间 $[1, 5]$ 内，没有任何正整数能同时承受 $2$ 与 $3$ 的整除，故一切共同正承载者必满足 $\text{主} \ge 6$；而 $6$ 本身既为正数，又兼备 $6 \bmod 2 = 0$ 与 $6 \bmod 3 = 0$ 之品格。既然``主''在一切兼具偶奇整除性的正数中居于最简、最先之必然下界，则必有 $\text{主} \le 6$。上下相合，遂得 $\text{主}=6$。

\medskip
\noindent\textbf{第四路：由完满性之差等推证首要完全数。}\\
诸数按其真部分和相对于自身的亏盈，呈现出完善性的高下差等：或因部分不足而亏缺，或因部分泛滥而过剩，唯有真部分和恰等于自身的完全数达致无亏无盈的均衡~\cite{aquinas,augustine}。依第四路完满居首规约，本原之至高地位不取决于数量之臃肿叠加，而在于位居完满序列之最先下界。由第一路对 $1$ 至 $5$ 的逐项排查已知，小于 $6$ 的自然数无一完满，故一切完全数皆满足 $\text{主} \ge 6$；又因 $6$ 自身已具足完全数之品格，任何位居全体完全数之首的本原必满足 $\text{主} \le 6$。由此，首要完全数唯独归于 $\text{主}=6$。

\medskip
\noindent\textbf{第五路：由部分配额之合目的秩序推证终极治理者。}\\
凡治理一整体者，必按恰当比例将份额授予诸部分，使各部分各得其所而又共同归于整体之目的~\cite{aquinas,metaphysics}。今考察如下配额秩序：治理者 $\text{主} > 0$ 先以二元与三元整除法则（$2 \mid \text{主} \land 3 \mid \text{主}$）确保二分之半（$\text{主}/2$）与三分之份（$\text{主}/3$）皆为完整不碎的真整数份额，绝无余数遗落；进而令半数份额、三分之一份额与单一主权单位（$1$）相加，恰好闭合于整体自身（$\text{主}/2 + \text{主}/3 + 1 = \text{主}$）。在整除约束下，半数与三分之一合占整体之六分之五，所余六分之一恰为单一单位 $1$，故必有 $\text{主}=6$。在公开源码 \texttt{lean/Main.lean} 中，五路主定理及联合充要会归定理进一步确立五路规约在全体自然数上均与 $n=6$ 互为充要条件；全部 20 条定理（含五路主定理、充要会归定理、共同存在性定理及异议辨析定理）均通过 Lean 4.34.1 内核验证。

\subsection{候选数之真因子和与本原资格核验}

表~\ref{tab:candidates} 对照列出了各典型候选数在五路标准下的真因子、真因子和及本原资格。

\begin{table}[htbp]
  \centering
  \caption{典型候选数的真因子、真因子和及本原资格}
  \label{tab:candidates}
  \small
  \begin{tabularx}{\textwidth}{c >{\raggedright\arraybackslash}p{3.4cm} c >{\raggedright\arraybackslash}X}
    \toprule
    候选数 $n$ & 真因子列表 $P(n)$ & 真因子和 $s(n)$ & 亏盈状态与本原资格 \\
    \midrule
    $0$ & \texttt{[]} & $0$ & 虽 $s(0)=0$，但非正数，不具现实存在资格 \\
    $1$ & \texttt{[]} & $0$ & 亏数（$s(1)<1$），无真部分，在流溢映射下由 $1\mapsto 0$ 失位 \\
    $2, 3, 5, 7$ & \texttt{[1]} & $1$ & 素数皆亏（$s(p)=1<p$），仅具单一真部分，不达自身总体 \\
    $4$ & \texttt{[1, 2]} & $3$ & 亏数（$s(4)=3<4$）；虽截断除法下 $4/2+4/3+1=4$，因 $3\nmid 4$ 失格 \\
    $6$ & \texttt{[1, 2, 3]} & $6$ & 完全数且居首，唯一同时满足五路规定之``主'' \\
    $10$ & \texttt{[1, 2, 5]} & $8$ & 亏数（$s(10)=8<10$），如奥古斯丁所记其真部分之和仅为 $8$~\cite{augustine} \\
    $12$ & \texttt{[1, 2, 3, 4, 6]} & $16$ & 盈数（$s(12)=16>12$），真部分之和逾越整体~\cite{augustine} \\
    $28$ & \texttt{[1, 2, 4, 7, 14]} & $28$ & 完全数（$s(28)=28$）~\cite{euclid}，然因 $6<28$ 失去首要本原资格 \\
    \bottomrule
  \end{tabularx}
  \par\vspace{0.4em}
  {\footnotesize\raggedright 数据来源：本文定义与形式验证。\par}
\end{table}

\section{讨论：经院异议之陈列与答复}

依循阿奎那《神学大全》``先立异议、再作答复''的辩难轨范~\cite{aquinas}，本节针对``$\text{主}=6$''这一命题所可能遭遇的四项诘难逐一予以剖析。

\subsection{异议一：关于``零''之不动性与``一''之至简性}

\noindent\textbf{异议}：似宜认``主''为 $0$ 或 $1$。盖第一路以 $s(n)=n$ 为不动之征，而 $0$ 满足 $s(0)=0$，岂非先于 $6$ 之不动者？又单子``一''纯一不可分，向称万数之源，岂非比复合之 $6$ 更配称``主''？

\noindent\textbf{答复}：如阿奎那所辨，第一本原乃完全之现实而非虚无~\cite{aquinas}。数 $0$ 之 $s(0)=0$ 仅系因空得零的消极静止，故诸路均以正性条件 $0<n$ 斥之。至于单子 $1$，因无真部分，在流溢下即跌落为 $0$，既不能自守定在，亦无从在第三路中整除承载偶与奇；故 $1$ 仅系 $6$ 的首个真部分，唯有统摄 $1, 2, 3$ 的 $6$ 方为现实自足之不动本原。

\subsection{异议二：关于次级完全数``二十八''之同等完满性}

\noindent\textbf{异议}：纵使``主''必为完全数，据欧几里得命题三十六，$28$ 之真因子和亦等于自身~\cite{euclid}；既同具完满性且 $28>6$，按数大为尊之常情，岂非 $28$ 更当居``主''位？

\noindent\textbf{答复}：此说混淆了数量多寡与本原先后。奥古斯丁指明六为首个完全数~\cite{augustine}；依第四路居首规约，若强立 $28$ 为首，必致 $28\le 6$ 之荒谬。况且 $28$ 含五项真因子且不被 $3$ 整除，皆违第二、五路规定，故仅为分受完满性的次级后裔，唯 $6$ 独居五路之首。

\subsection{异议三：关于自然数截断除法下``四''对第五路之伪冒}

\noindent\textbf{异议}：在自然数截断除法下，$n=4$ 亦满足 $4/2+4/3+1=4$，且 $4<6$，何以断言第五路唯独指向 $6$？

\noindent\textbf{答复}：分整体为三而弃余数，乃配额之折损而非完满。数 $4$ 作三分时暗弃余数 $1$，仅借向下取整掩盖亏空；故第五路先行确立整除约束 $2\mid n\land 3\mid n$，将违背模 $3$ 整除的 $4$ 排除在外。唯有 $\text{主}=6$，方能以真正的半数（$3$）、三分之一（$2$）与单位（$1$）无余闭合。

\subsection{异议四：关于五路规定之交织与同一本原}

\noindent\textbf{异议}：第一与第四路皆赖小于 $6$ 无完全数之事实，第三与第五路亦共含 $2$ 与 $3$ 整除条件，五路岂非并非彼此绝缘之系统？

\noindent\textbf{答复}：阿奎那立五路，本是从受造秩序的五个不同起点推证同一本原，而非判归五个相异实体~\cite{aquinas}。本文五项规约各据独立设问立论，其底层算术性质彼此呼应，恰印证诸路在数之秩序中终必会归同一本原 $\text{主}=6$。

详辨之，第一路旨在排除较小正不动点，第四路则对全体完全数求最小，二者展开 $K$ 后等价，并非独立公理；第三路依通常自然数序求 2、3 最小正公倍数，第五路则为整除约束下之配额方程。

\section{结论}

本文依循阿奎那``五路证明''体制与古典完全数传统，在 Lean 4（4.34.1）中将第一本原``主''设为待定自然数参数，在不预设常量与新增公理下，证明流溢静止、原初三因、偶奇承载、完满居首及治理配额五路皆充要导出 $\text{主}=6$，并排除 $0, 1, 10, 12, 28$ 与 $4$ 诸数。故``$\text{主}=6$''实乃算术秩序在五重维度下不可移易之归宿。

\section*{声明}

\subsection*{数据与代码可用性声明}
本文 Lean 4 源码（\texttt{lean/Main.lean}）、核验记录（\texttt{evidence/verification.json}）及复现说明（\texttt{README.md}）已公开于 GitHub 仓库 \url{https://github.com/chengxi271-commits/zhu-six-five-ways}（首个提交 \texttt{2aa62b8}），可在标准 Lean 4.34.1 环境下直接复现。

复现时自仓库 \texttt{lean} 目录执行 \texttt{lake build}，再返仓库根目录运行 \texttt{python} \texttt{scripts/verify.py} \texttt{--lean lean}，即生成含 \texttt{\#print axioms} 之逐定理审计：20 项定理仅用标准公理，无 \texttt{sorryAx}。

\subsection*{作者贡献声明}
Chan Isidore负责本研究的概念化、五路论证方向确立、经院体例设计、文稿审阅修订与整体监督。

\subsection*{利益冲突声明}
作者声明无利益冲突。

\subsection*{致谢}
在本研究开展过程中，Codex 用于 Lean 4 形式化与验证，Gemini 用于论文文稿写作；作者对全文内容承担全部责任。

\begin{thebibliography}{99}

\bibitem{aquinas}
AQUINAS T. Summa theologiae: I, q. 2, aa. 2-3[M/OL]. Fathers of the English Dominican Province, trans. 2nd rev. ed. New Advent, 1920[2026-11-14]. \url{https://www.newadvent.org/summa/1002.htm}.

\bibitem{augustine}
AUGUSTINE. The city of God: XI, 30[M/OL]. DODS M, trans. Buffalo: Christian Literature Publishing Co., 1887[2026-11-14]. \url{https://www.newadvent.org/fathers/120111.htm}.

\bibitem{euclid}
EUCLID. Elements: Book IX, proposition 36[M/OL]. JOYCE D E, ed. Worcester: Clark University, 1996[2026-11-14]. \url{https://mathcs.clarku.edu/~djoyce/elements/bookIX/propIX36.html}.

\bibitem{physics}
ARISTOTLE. Physics: Book VIII[M/OL]. HARDIE R P, GAYE R K, trans. Cambridge, MA: The Internet Classics Archive, Massachusetts Institute of Technology[2026-11-14]. \url{https://classics.mit.edu/Aristotle/physics.8.viii.html}.

\bibitem{metaphysics}
ARISTOTLE. Metaphysics: Book II[M/OL]. ROSS W D, trans. Cambridge, MA: The Internet Classics Archive, Massachusetts Institute of Technology[2026-11-14]. \url{https://classics.mit.edu/Aristotle/metaphysics.2.ii.html}.

\bibitem{lean2015}
DE MOURA L, KONG S, AVIGAD J, et al. The Lean theorem prover (system description)[C]//Automated Deduction---CADE-25. Cham: Springer, 2015: 378-388. DOI: 10.1007/978-3-319-21401-6\_26.

\bibitem{lean4}
DE MOURA L, ULLRICH S. The Lean 4 theorem prover and programming language[C]//Automated Deduction---CADE 28. Cham: Springer, 2021: 625-635. DOI: 10.1007/978-3-030-79876-5\_37.

\bibitem{tpil}
AVIGAD J, DE MOURA L, KONG S, et al. Theorem proving in Lean 4[M/OL]. Lean community[2026-11-14]. \url{https://lean-lang.org/theorem_proving_in_lean4/}.

\bibitem{modal}
BENZMÜLLER C, WOLTZENLOGEL PALEO B. The Inconsistency in Gödel’s Ontological Argument: A Success Story for AI in Metaphysics[C]//IJCAI. 2016: 936–942.
\end{thebibliography}

\end{document}
